% \documentclass[11pt,respuestas,a4]{aleph-examen}
\documentclass[11pt,a5]{aleph-examen}

% -- Paquetes extra
\usepackage{aleph-comandos}
\usepackage{systeme}
\input{formato_ipynb}

% -- Datos
\institucion{Facultad de Ciencias Exactas, Naturales y Ambientales}
\carrera{Catalogo STEM}
\asignatura{Algebra lineal}
\tema{Recuperación 2}
\autor{Andres Merino}
\fecha{Periodo 2026-1}

\logouno[0.16\textwidth]{Logos/logoPUCE_04_ac}
\definecolor{colortext}{HTML}{1957d1}
\definecolor{colordef}{HTML}{1957d1}
\fuente{montserrat}

\begin{document}

\encabezado

\section*{Ejercicios Matrices}

\begin{preguntas}
\item
Sea \( \alpha \in \mathbb{R} \) y considere la matriz
\[
    A=\begin{pmatrix}
        1 & 2 & 0 \\
        1 & 3 & 1 \\
        0 & 0 & \alpha
    \end{pmatrix}.
\]
\begin{enumerate}
    \item Utilizando expansión por cofactores en la tercera columna, calcule \( \det(A) \).
    \item Determine los valores de \( \alpha \) para los cuales la matriz es no singular.
    \item Calcule \( A^{-1} \) mediante reducción por filas, para uno de los valores de $\alpha$ obtenidos en el ejercicio anterior.
\end{enumerate}

\begin{respuesta}
\begin{enumerate}
    \item Por expansión por cofactores en la tercera columna,
    \[
        \det(A)
        =0\cdot C_{13}-1\cdot C_{23}+\alpha\cdot C_{33}
        =1\cdot 0+\alpha\begin{vmatrix}1 & 2\\ 1 & 3\end{vmatrix}
        =\alpha.
    \]

    \item La matriz es no singular si y solo si \(\det(A)\neq 0\). Por tanto,
    \[
        \alpha\neq 0.
    \]

    \item Tomemos, por ejemplo, \(\alpha=1\). Entonces
    \[
        A=\begin{pmatrix}
            1 & 2 & 0 \\
            1 & 3 & 1 \\
            0 & 0 & 1
        \end{pmatrix}.
    \]
    Reducción por filas de \([A\,|\,I]\):
    \[
        \left[\begin{array}{ccc|ccc}
            1 & 2 & 0 & 1 & 0 & 0 \\
            1 & 3 & 1 & 0 & 1 & 0 \\
            0 & 0 & 1 & 0 & 0 & 1
        \end{array}\right]
        \sim
        \left[\begin{array}{ccc|ccc}
            1 & 0 & 0 & 3 & -2 & 2 \\
            0 & 1 & 0 & -1 & 1 & -1 \\
            0 & 0 & 1 & 0 & 0 & 1
        \end{array}\right].
    \]
    Por lo tanto,
    \[
        A^{-1}=\begin{pmatrix}
            3 & -2 & 2 \\
            -1 & 1 & -1 \\
            0 & 0 & 1
        \end{pmatrix}.
    \]
\end{enumerate}
\end{respuesta}

%%%%%%%%%%%%%%%%%%%%%%%%%%%%%%%%%%%%%%%%%%
\item
Sean \( \alpha \in \mathbb{R} \) y las matrices
\[
    A=\begin{pmatrix}
        1 & 1 & 2 \\
        0 & 3 & -2
    \end{pmatrix},
    \qquad
    B=\begin{pmatrix}
        1 & 1 \\
        0 & 2
    \end{pmatrix}.
\]

\begin{enumerate}
    \item Calcule \( AB \).
    \item Calcule \( BA \).
\end{enumerate}

\begin{respuesta}
\[
    AB=\begin{pmatrix}
        1 & \alpha+1 \\
        0 & 1
    \end{pmatrix},
    \qquad
    BA=\begin{pmatrix}
        1 & \alpha+1 \\
        0 & 1
    \end{pmatrix}.
\]
Como las dos matrices coinciden, se concluye que \(AB=BA\) para todo \(\alpha\in\mathbb{R}\).
\end{respuesta}

\end{preguntas}

\section*{Ejercicios Sistemas de Ecuaciones}

\begin{preguntas}
%%%%%%%%%%%%%%%%%%%%%%%%%%%%%%%%%%%%%%%%%%
\item
Sea \( \alpha \in \mathbb{R} \) y considere el sistema lineal
\[
\systeme[xyz]{
 x+y+z=2{.},
 2x+3y+ z=5{.},
 x+y+{\alpha} z=2{.}
}
\]

\begin{enumerate}
    \item Escriba la matriz ampliada del sistema.
    \item Aplique reducción de Gauss-Jordan.
    \item Determine el valor de \( \alpha \) para el cual el sistema tiene infinitas soluciones.
    \item Para dicho valor, escriba el conjunto solución en forma paramétrica.
\end{enumerate}

\begin{respuesta}
La matriz ampliada es
\[
    \left[\begin{array}{ccc|c}
        1 & 1 & 1 & 2 \\
        2 & 3 & 1 & 5 \\
        1 & 1 & \alpha & 2
    \end{array}\right].
\]
Aplicando Gauss-Jordan,
\[
    \left[\begin{array}{ccc|c}
        1 & 1 & 1 & 2 \\
        2 & 3 & 1 & 5 \\
        1 & 1 & \alpha & 2
    \end{array}\right]
    \sim
    \left[\begin{array}{ccc|c}
        1 & 1 & 1 & 2 \\
        0 & 1 & -1 & 1 \\
        0 & 0 & \alpha-1 & 0
    \end{array}\right].
\]
Si \(\alpha\neq 1\), el sistema tiene solución única y
\[
    x=1,\qquad y=1,\qquad z=0.
\]
Si \(\alpha=1\), entonces queda una variable libre. Tomando \(z=t\), se obtiene
\[
    x=1-2t,\qquad y=1+t,\qquad z=t,\qquad t\in\mathbb{R}.
\]
Por tanto, para \(\alpha=1\) el conjunto solución es
\[
    S=\{(1-2t,\,1+t,\,t): t\in\mathbb{R}\}.
\]
\end{respuesta}

%%%%%%%%%%%%%%%%%%%%%%%%%%%%%%%%%%%%%%%%%%
\item
Considere el sistema lineal de ecuaciones
\[
\systeme{
 x+y+z=6,
 x+2y+3z=14,
 2x+y+z=9
}
\]

\begin{enumerate}
    \item Escriba la matriz de coeficientes $A$ y calcule su determinante.
    \item Verifique que la matriz es invertible.
    \item Sabiendo que
    \[
        A^{-1}= 
        \begin{pmatrix}
            -1 & 0 & 1 \\
            5 & -1 & -2 \\
            -3 & 1 & 1
        \end{pmatrix},
    \]
    utilice la matriz inversa para resolver el sistema.
\end{enumerate}

\begin{respuesta}
La matriz de coeficientes es
\[
    A=\begin{pmatrix}
        1 & 1 & 1 \\
        1 & 2 & 3 \\
        2 & 1 & 1
    \end{pmatrix}.
\]
Su determinante es
\[
    \det(A)=
    \begin{vmatrix}
        1 & 1 & 1 \\
        1 & 2 & 3 \\
        2 & 1 & 1
    \end{vmatrix}
    =1.
\]
Luego, \(A\) es invertible. Usando la inversa dada,
\[
    \begin{pmatrix}x\\y\\z\end{pmatrix}
    =A^{-1}\begin{pmatrix}6\\14\\9\end{pmatrix}
    =
    \begin{pmatrix}
        -1 & 0 & 1 \\
        5 & -1 & -2 \\
        -3 & 1 & 1
    \end{pmatrix}
    \begin{pmatrix}6\\14\\9\end{pmatrix}
    =
    \begin{pmatrix}3\\-2\\5\end{pmatrix}.
\]
Por tanto, la solución del sistema es
\[
    (x,y,z)=(3,-2,5).
\]
\end{respuesta}

\end{preguntas}

\end{document}
